\chapter*{Summary}
\addcontentsline{toc}{chapter}{Summary}

!! THIS IS AN AI GENERATED SUMMARY. PLEASE REVIEW AND EDIT AS NEEDED. !!

This thesis focused on the development of vehicle speed estimation algorithms based on artificial intelligence methods. The main objective was to develop a direct vehicle speed estimator using recurrent neural networks (RNNs) and compare its performance to a conventional wheel based speed estimator enhanced with AI based components. The direct speed estimator was developed using a dataset recorded from a test vehicle equipped with various sensors. The dataset included measurements such as wheel speeds, accelerations, yaw rate, and the ground truth vehicle speed obtained from a high-precision GPS system. The RNN architecture was chosen due to its ability to capture temporal dependencies in sequential data, making it suitable for vehicle speed estimation tasks. The development process involved data preprocessing, feature selection, model training, and evaluation. Various RNN architectures were explored, including Long Short-Term Memory (LSTM) and Gated Recurrent Unit (GRU) networks. Hyperparameter tuning was performed to optimize the model's performance. The direct speed estimator was evaluated using metrics such as root mean squared error (RMSE) and percentage within threshold (PWT). The results demonstrated that the RNN-based estimator achieved high accuracy in estimating vehicle speed, outperforming traditional methods in certain scenarios.

In addition to the direct speed estimator, a conventional wheel based speed estimator was developed and enhanced with AI based components. This estimator utilized wheel speed measurements and incorporated machine learning techniques to improve its accuracy. The performance of the conventional estimator was compared to the RNN-based direct estimator using the same evaluation metrics. The comparison revealed that while the conventional estimator performed well under normal driving conditions, the RNN-based estimator exhibited superior performance in scenarios involving rapid acceleration or deceleration, as well as varying road conditions.

Overall, this thesis demonstrated the effectiveness of RNNs in vehicle speed estimation tasks and highlighted the potential of AI based methods to enhance traditional estimation techniques. The developed direct speed estimator showed promising results and could be further improved with additional data and advanced architectures. Future work could explore the integration of other sensor modalities, such as camera or LiDAR data, to further enhance the accuracy and robustness of vehicle speed estimation algorithms.

